\chapter*{Class and Party}
\addcontentsline{toc}{chapter}{Class and Party}
\setcounter{postnote}{0} % Restarting endnotes numbering for each chapter

\section*{1}

Ruling classes have at their disposal a formidable range of weapons for the maintenance of their domination and the defence of their power and privileges. How then are these ruling classes to be undone, and how is a new social order to be established?

That this can be achieved is an essential theme of the Marxist message: but a no less essential theme of that message is that it does have to be achieved. No doubt, the desired transformation must largely depend upon, or at least be linked with, the deepen- ing of the contradictions of capitalism and their consequent and manifold impact upon the superstructure. But even so, the transformation will have to be brought about by human intervention and practice, and will be the result of growing class conflict and confrontation, in which the working class must play a predominant role.

In order that it may play that role effectively, there must be organization: against the vast array of powerful forces which the ruling class is able to deploy in the waging of class struggle, the working class and its allies cannot hope to succeed unless they are organized.

The question is what this means; and there have on this score been very wide divergences within the spectrum of Marxist political thought. These divergences must not be translated into a simplistic and misleading contraposition of 'organization' and 'spontaneity'. There is no Marxist thinker, of any sort, who has ever advocated pure spontaneity as a way of revolutionary practice. The notion is evidently absurd: even a levee en masse requires to be organized if it is to get anywhere.

Nor, at the other extreme, has anyone within the Marxist tradition ever advocated that a revolution should be made by an organized group or party without any measure of popular support. This has at one time or other been put forward as a means of revolutionary action and change; but it does not form part of any variety of Marxism, and has in fact been consistently rejected and denounced by Marxists as mere 'putschism'.

The question then is not, in Marxist terms, of a 'choice' between organization and spontaneity, or between the party and the class. There is in reality no such choice; and Marxists have quite rightly insisted that this was for them a false dichotomy. But the divergences have nevertheless been very real, and have centred round different views of the relation of the working class to the organization, and of the relative weight to be attributed to either. The differences may be matters of emphasis; but emphasis in this realm is not a matter of detail.

These divergences point to a permanent tension in Marxist thought and practice in regard to the relation of class and party, not surprisingly since the issues which that relation raises are crucial, not only to the advancement of the revolutionary process but to the shape and character of the Marxist project itself. The attempt has often been made to blur and negate this tension by reference, here as in other questions, to the 'dialectical' interrelation which must be established between class and party. This is no doubt desirable, but to say so is not to resolve anything; it merely restates the problem in different words.

Given the wide spectrum of views and attitudes on the question of class and party, it is as well to note first that Marx himself stands at one end of the spectrum, in so far as his own emphasis falls very heavily on the action of the class. His concern throughout his political life was not simply with the emancipation of the working class: this could be said to be the avowed purpose of all revolutionaries. His concern is with the emancipation of the working class by its own efforts. He expressed this most suc- cintly in the Preamble to the Provisional Rules of the First International in 1864. 'That the emancipation of the working classes must be conquered by the working classes themselves'.\postnote{FI, p. 82.}

Neither this nor anything else Marx said on the subject precludes organization in the form of a party of the working class; and there are in his work innumerable references to the need for the working class to organize itself. On the other hand, he was not particularly concerned with the form which the political organization of the proletariat should assume, and was content to leave workers in different countries to determine this according to their own circumstances.\postnote{See, e.g., Marx's insistence on this point in an interview in The World in July 1871, in Werke (Berlin, 1964), op. cit. XVII, p. 641. See also M. Johnstone, 'Marx and Engels and the Concept of the Party' in The Socialist Register 1967 (London, 1967).} The one definite prescription was that the party should not be a sect, isolated from the working class, and composed of the 'professional conspirators' whom he savagely denounced as the 'alchemists of revolution' in 1850.\footnote{'It goes without saying', Marx wrote ironically, 'that it is not enough for them to engage in the organisation of the revolutionary proletariat. Their business consists precisely in the anticipation of revolutionary development, in artificially bringing it to a crisis, in making a revolution by improvisation without the conditions of a revolution. For them, the only condition for a revolution is the proper organisation of their conspiracy. They are the alchemists of revolution and are quite like the earlier alchemists in their narrowness of mind and fixed prejudices' (Werke, (Berlin, 1960), op. cit. VII, p. 273).}

Whatever the form of the party, it is the working class, its developing consciousness and its struggles for self-emancipation which really matter to Marx: the party is only the political expression and the instrument of the class. Even this formulation may go further than is warranted in Marx's case and it may be more accurate to interpret him, in this connection, as seeing the working class itself performing its political role, with the political party helping it to do so. The formulation is rather ambiguous—but so are Marx's pronouncements on this issue.

What is not at all ambiguous is the faith which Marx and Engels had in the capacity of the working class to achieve selfemancipation. In a Circular Letter to the leadership of the German Social-Democratic Workers' Party, written in 1879, they angrily rejected any flirtation with the idea that 'the working class is incapable of liberating itself by its own efforts' and that 'for this purpose it must first accept the leadership of “educated and propertied” bourgeois, who also have “opportunity and time ' ' to acquaint themselves with what is good for the workers.'\postnote{FI, p. 370.} They reminded the leaders of the party that 'when the International was formed, we expressly formulated the battle- cry: the emancipation of the working class must be the work of the working class itself', and that 'we cannot ally ourselves, therefore, with people who openly declare that the workers are too uneducated to free themselves and must first be liberated from above by philanthropic big bourgeois and petty bourgeois.'\postnote{Ibid., p. 375.} The point they were making is clearly of wider application; what they were concerned to stress was their general belief in the self-emancipating capacities of the working class. The question here is not whether they were right or wrong, and what are the problems they were overlooking. The fact is that for them, the class always came first, the party a long way behind. This cuts very deep, and has a direct bearing on the wider question of the direct and indirect exercise of popular power, and on the meaning of socialist democracy.\postnote{See below, pp. 138 ff.}

By the time Engels died in 1895, German Social Democracy was well on the way to becoming an authentic mass organization, and it continued thereafter to make great progress. As one writer has noted, the German Social Democratic Party had by 1914 become a vast institution that was staffed by more than 4,000 paid functionaries and 11,000 salaried employees, had 20,000,000 marks invested in business and published over 4,000 periodicals'.\postnote{H. Gruber, International Communism in the Era of Lenin (New York 1972) p. 12.} It also had a very substantial parliamentary representation and was a force in provincial and local government.

To a greater or lesser extent, much the same was coming to be true of other European Social Democratic parties; and it was more or less taken for granted that the parties of the working class would be, and indeed must be, mass parties, deeply involved in the political life of their respective countries, though loosely linked in the Second International. With this growth of working-class parties, and the expectation of their further implantation and progress, there also went a very general, though not of course unanimous, acceptance that the transformation of capitalist society must be envisaged as a strictly constitutional process, which must on no account be endangered by an ill-conceived activism and adventurist policies. How widespread this acceptance was, inside German Social Democracy but also elsewhere, was somewhat masked by the opposition which Bernstein aroused by his explicit 'revisionism'. But in more accentuated forms or less, with variations that were less important than rhetoric was intended to suggest, 'revisionism' was the perspective which dominated all but a small part of European Social Democracy: the great 'betrayal' of 1914—and of the years of war—was a natural manifestation of it.

In the present context, what is important about this perspective is that it inevitably led to the exaltation of the party as the accredited and ever more influential representative in national life of the working class, as the expression of its political presence at all levels, but also as its guardian against those who pressed upon the working class actions and policies which the leadership of the party judged and proclaimed to be irresponsible and dangerous. The exaltation of the party also meant the enhancement of the status and position of the party leaders, the men who were in charge of that complex and delicate machinery whereby the locomotive of socialism was to be driven, at safe speed, through capitalist society.

This is not to endorse Robert Michels's 'iron law of oligarchy', which was formulated in his Political Parties (1915) with the German Social Democratic Party as a main point of reference.

But we must note here that there was, in the notion and practice of the working-class mass party, integrated in the constitutional framework of bourgeois society, a powerful logic driving towards the concentration of power in the hands of leaders claiming to be the representatives and spokesmen of the working class, and occupying a privileged position in the determination of policy. It was not Lenin who started it all.

Lenin had in fact an extraordinarily strong sense of the need both to build a party of a special kind, given the conditions of Czarist Russia, and to maintain as close a link as possible with the working class. Without the constant reinvigoration of the party by the working class, it would go stale and bureaucratic, and Bag behind' the masses—one of Lenin's constant fears. He had conceived from his earliest days in the revolutionary movement of a party of 'professional revolutionaries', adapted to the struggle against a despotic regime by a highly centralized command structure, which came to be called 'democratic centralism'. He did not suggest, before 1914, that the party he wanted for Russia was suitable for what he called 'countries where political liberty exists'.\postnote{V. I. Lenin, What is to be done? p. 106.} The important point for him indeed the essence of Lenin's contribution to Marxism was that there must be organization and direction if the revolutionary process was to be advanced.

Lenin did not fear the passivity of the working class, but that its struggles would lack political effectiveness and revolutionary purpose. It was to this end that the party was essential: without its guidance and leadership the working class would be a social force capable of spasmodic and incoherent actions, but incapable of turning itself into the disciplined army that was required to overthrow Czarism and to advance towards the conquest of socialist power.

Even so, it needs to be stressed that Lenin knew well that the party could not fulfil its tasks without being steeped and involved in the experiences of the masses; and whenever opportunity offered for the party to 'open up'—as in 1905 and 1917 and after—he seized the chance and mercilessly castigated the 'party bureaucrats' steeped in their routine.\postnote{See, e.g., Liebman, Leninism under Lenin.} Even in those of his writings which most insistently focus on the need for organization, centralism, discipline, etc., and of which What is to be Done? is an extreme example,\footnote{Lenin noted some years after the publication of his pamphlet that, in the circumstances in which it was written, he had had to 'bend the stick' somewhat, and that some of his formulations had been deliberately polemical. This however was no repudiation of the basic themes of the argument.} the need for an organic connection with the working class is never ignored. And it was also Lenin who wrote the following in 1920:
\begin{displayquote}
    History as a whole, and the history of revolutions in particular, is always richer in content, more varied, more multiform, more lively and ingenious than is imagined by even the best parties, the most classconscious vanguards of the most advanced classes. This can readily be understood, because even the finest of vanguards express the class- consciousness, will, passion and imagination of tens of thousands, whereas at moments of great upsurge and the exertion of all human capacities, revolutions are made by the class-consciousness, will, passion and imagination of tens of millions, spurred on by a most acute struggle of classes.\postnote{'Left-Wing' Communism—An Infantile Disorder, in SWL, p. 574.}
\end{displayquote}

This, however, did not in the least lead him to decry the crucial importance of the party in the revolutionary process, and to devalue its role in relation to the working class. He could never have said, as Rosa Luxemburg said in 1904, that 'historically, the errors committed by a truly revolutionary movement are infinitely more fruitful than the infallibility of the cleverest Central Committee.'\postnote{R. Luxemburg, Organizational Question of Social Democracy, in Rosa Luxemburg Speaks, p. 130.} He did not believe in the 'infallibility' of any Central Committee or party organ. But neither did he believe that a 'truly revolutionary movement' was conceivable that did not include a well-structured organization. It was in the same work from which the earlier quotation is taken that Lenin observed, in answer to those in the German Communist movement who counterposed 'dictatorship of the party or dictatorship of the class; dictatorship (party) of the leaders, or dictatorship (party) of the masses', that '. . . as a rule and in most cases—at least in present-day civilized countries—classes are led by political parties . . . political parties, as a general rule, are run by more or less stable groups composed of the most authoritative, influential and experienced members, who are elected to the most responsible positions, and are called leaders'.\postnote{Lenin, op. cit., p. 532.}

In saying this, Lenin was pointing to the obvious but essential (and to many the rather uncomfortable and disagreeable) fact that any 'model' of the revolutionary process, or for that matter of working-class politics in general, would have to include an organized political formation, which also meant that there would be a leadership and a structure of command. Lenin's whole emphasis in this respect was different from that of Marx and Engels. As was noted earlier, they too had a concept of the party: yet the shift of emphasis is unmistakable. But this, to a greater or lesser extent, was true of all strands of Marxist thought: well before 1914, the international Marxist movement, West as well as East, had developed 'models' of the party, both of which attributed much greater importance to the party than Marx and Engels had ever thought necessary, though the point is less certain in relation to Engels in his last years, when he was closely involved in the affairs of the German Social Democratic Party.

It was not of course on the issue of the party itself that Lenin's opponents attacked him, but on the dictatorial centralism which they accused him of advocating, and for seeking, as they suspected and alleged, his own personal dictatorship through the centralized and subordinate party. In 1904, after the great split of 1903 between Bolsheviks and Mensheviks and the subsequent publication of Lenin's ultra-centralist One Step Forward, Two Steps Back, Trotsky published Our Political Tasks, in which he warned against the dangers of 'substitutism' and issued the prophecy to which much later events were to give an exceptional resonance: Lenin's methods lead to this: the party organization at first substitutes itself for the party as a whole; then the Central Committee substitutes itself for the organisation; and finally a single “dictator” substitutes himself for the Central Committee.\postnote{I. Deutscher, The Prophet Armed. Trotsky: 1879-1921 (London 1954) p. 90.} Like Lenin s Menshevik opponents, he wanted a 'broadly based party' and insisted that 'the party must seek the guarantee of its stability in its own base, in an active and self- reliant proletariat, and not in its top caucus\dots{}'\postnote{Ibid., p. 90.}

There was a great deal in this, as in the whole debate, which was rhetorical and illusory, in so far as many of the participants tended to exaggerate very considerably what really separated them from Lenin and the 'centralists'. This was certainly true of Trotsky, who was in his years of power to exhibit even more 'centralist' tendencies than Lenin had done or was later to do. Lenin would hardly have demurred from the view that there was need for an 'active' proletariat: he never ceased to make the point himself. But what did 'self-reliant' exactly mean? Surely not, for  Trotsky, that the proletariat could do without the party. As for the need to have a 'broadly based' party, this too would depend on the meaning of the concept: as Lenin showed within a year of the publication of Trotsky's pamphlet, that is to say in the course of the Russian Revolution of 1905, he had no wish whatever to create a party that was small, self-contained and inward- looking. When 1905 appeared to open new possibilities, he seized them eagerly; and so too in 1917.

In her own contribution to the debate after the split of 1903, namely Organizational Question of Social Democracy, which was also published in 1904, Rosa Luxemburg attacked Lenin for his 'ultra-centralism' and argued that 'social democratic centralism cannot be based on the mechanical subordination and blind obedience of the party membership to the leading party centre'.\postnote{R. Luxemburg, Organizational Question of Social Democracy, in Rosa Luxemburg Speaks, p. 118.} Lenin would of course have denied that this is what he wanted. But more important than the accusation was Rosa Luxemburg's own acceptance of a certain kind of centralism—what she called 'social democratic centralism', which she opposed to 'unqualified centralism'. But save for her insistence that the party must on no account stifle the activity of the movement, she was unable to define what her centralism would entail, particularly for a regime like the Russian, and she indeed noted that 'the general ideas we have presented on the question of socialist centralism are not by themselves sufficient for the formulation of a constitutional plan suiting the Russian party';\postnote{Ibid., p. 122.} 'In the last instance', she added, 'a statute of this kind can only be determined by the conditions under which the activity of the organization takes place in a given epoch.'\footnote{Ibid., p. 122. It may also be noted that her celebration of the mass strike, based on the experience of the Russian Revolution of 1905, was intended to incorporate this means of struggle into the political strategy of Marxism, and to dissociate it from anarchism. See R. Luxemburg, The Mass Strike, The Political Party and the Trade Unions, ibid., p. 158.} At the time, Lenin would hardly have disagreed with that.

However, there was much more to the debate on party and class than the question of centralism. What was in fact at stake was the very nature of Marxist politics in relation to the class or classes whose emancipation was its whole purpose. The issues involved acquired an entirely new dimension with the Bolshevik seizure of power in 1917; and they have remained central to Marxist politics right to the present day.

\section*{2}
The first of these issues concerns the representativeness of working-class parties. The early debate which opposed Luxemburg to Lenin had largely turned on the role of the party in relation to the working class. Lenin did not believe that the latter would, without proper leadership, become a truly revolutionary force, while Luxemburg feared that the party, if it came to control the working class, would stifle its militant and creative impulses.\footnote{Ibid., p. 121. ' The tendency is for the directing organs of the socialist party to play a conservative role.'} Both, however, assumed—and so did everybody else who took part in the debate—that given the right form of organization and sufficient time, there was a point in the spectrum where class and party could form a genuinely harmonious and organic unity, with the party as the true expression of a classconscious and revolutionary working class. People might differ greatly over the question of how this could be achieved, and where the point on the spectrum was located: but they did not differ in the belief that there was such a point, and that it was their business to get to it. In other words, 'substitutism' was a danger, and might become a reality. But it was in no way ineluctable.

This failed to take into account the fact that some degree of substitutism is bound to form part of any kind of representative organization and of representative politics at all levels. Of course, tne degree matters very greatly, and can be very considerably affected, one way or the other, by various means and devices as well as by objective circumstances. But the notion of the party achieving an organic and perfectly harmonious representation of the class is nothing but a more or less edifying myth. Marx mostly avoided the problem by focusing on the class rather than the party, but the problem is there all the same. It is there, it may be noted, for the best of reasons, namely that Marxists have a commitment to thorough political democratization and to what may be called the dis-alienation of politics. Paities and movements which have no such commitment, or for whom it has much less importance, naturally have in this respect an easier time of it.

The problem of 'substitutism' arises in various forms and degrees because the working class, as was noted and discussed in Chapter II, is not a homogeneous entity, and that the 'unity of the working class', which the party seeks or claims to embody, must be taken as an exceedingly dubious notion, which may well come to have some definite meaning in very special and unusual circumstances, but which normally obscures the permanent and intractable differences and divisions which exist in this as in any other social aggregate.

On this view, a united party of the working class, speaking with one voice, must be a distorting mirror of the class; and the greater the 'unity', the greater the distortion, which reaches its extreme form in the 'monolithic' party.

Yet, political parties of the working class are not debating societies, to use a consecrated formula; and they do need some degree of unity in what is a permanent and often bitter class struggle. A genuinely 'representative' party, in which all the divisions of the working class find full expression by way of factions and tendencies and endless debate, may thereby find itself unable to cope with the very responsibilities which are the reason for its existence; and this is likely to be exceptionally true in periods of acute conflict and crisis. The demands of representatives on the one hand, and of effectiveness on the other, are not altogether irreconcilable, in that a more representative party may be more effective than one which lives by imposed and spurious 'unity'. But it is an illusion to think that the contradiction is not a genuine one. It is genuine; and although it may be capable of attenuation by democratic compromise and toleration, it is very unlikely to be fully resolved in any foreseeable future.

In any case, it is unrealistic to speak of 'the party' as if it could be taken for granted that there was one natural political organ of the working class, with the unique mission to represent it politically (and for that matter in many other ways as well).

It was perhaps inevitable that the concept of 'the party' should have come to be enshrined or at least accepted very early in the Marxist perspective of working-class politics. It will be recalled that in the Communist Manifesto, Marx and Engels claimed as the distinctive mark of 'the Communists' that 'in the national struggles of the proletarians of the different countries, they point out and bring to the front the common interests of the entire proletariat, independently of all nationality'; and that 'they always and everywhere represent the interests of the movement as a whole'. This being the case, they added, the Communists were 'practically, the most advanced and resolute section of the working-class parties of every country'; and they were also those who had the advantage over the great mass of the proletariat 'of clearly understanding the line of march, the conditions, and the ultimate general results of the proletarian movement'.\postnote{Revs., pp. 79-80.} Although these formulations left open the question of the form of organization which the Communists should adopt\footnote{The Communists do not form a separate party opposed to other working class parties', ibid., p. 79. The formulations which appear in the Manifesto are in this context by no means unambiguous—nor of course do they represent the authors' last word on the subject.} it is obvious that Marx and Engels were here talking of a definite vanguard; and the notion of more than one vanguard is absurd.

From the notion of a vanguard to that of a vanguard party, there was only a short step, which Russian revolutionaries in particular, given their specific circumstances, found it easy to take, and which they had very little option but to take. Moreover, the solidification of Marxism into 'scientific socialism' further helped to strengthen the view that there could only be one true party of the working class, namely the one which was the bearer of the 'truth' of the movement. This might not prevent the existence of other parties, but it made much more difficult the acceptance of the notion that there might be more than one Marxist working-class party. Indeed, it made that notion virtually unacceptable.

From a very different political tradition, German Social Democracy, as it developed in the last decades of the nineteenth century, was also driven towards the concept of the unified party of the working class, which would bring together not only different sections of the working class into a coherent political formation, but which would also be capable of bringing under its own banner other classes and strata, and thus become the party of all those who wanted a radical reorganization of the social order.

Nor, given the relative looseness in ideology and organization of this mass party, and the fact that it did attract a very large measure of working-class support, was there much encouragement, for those who found the party inadequate or worse, to cut loose and form their own party. It is highly significant in this connection that the German Social Democratic Party should have held together until well into the First World War, despite the intense alienation felt by a substantial fraction of its left- wing members. Even Rosa Luxemburg took it that 'the party', i.e. the existing German Social Democratic Party, was her only possible political home, notwithstanding her disagreements with it. By the time she and others changed their minds, or had their minds changed tor them by the circumstances of war, it was too late: the massive bulwark of the existing social order which the German Social Democratic Party effectively represented was intact, and the leaders of the party were both able and ready to sustain that social order in its hour of greatest need, in the winter of 1918/19.

The encouragement to retain a unified party is further strengthened by the hope of changing its policies; by the fear that splitting off must mean isolation, marginality, and ineffectiveness; and by even greater fear of acting as a divisive agent and thereby (so it is thought) weakening the working-class movement. These are very strong pressures, which clearly work in favour of the leadership and apparatus of the mass party, and strengthen their appeal to unite, close ranks, show loyalty and so forth.

Even so, the working-class movement in capitalist countries has not normally found its political expression in one single party. In some cases, one party has been able to establish a quasi-monopoly as the party of organized labour—the Labour Party in Britain being a case in point. But even in such cases, the labour movement has produced more than one party; and the quasi-monolistic party has had to fight very hard to maintain its position. It is clear that more-than-one-party is in fact the 'natural' expression of the politics of labour. 'The party' as the single legitimate expression of the labour movement is an invention which postdates the Bolshevik Revolution. There is nothing in classical Marxism which stipulates such singularity.

Given the heterogeneity of the working class and of the working-class movement, it would be very remarkable if one party did constitute its natural expression; and the point is reinforced by the extension of the notion of the working class which is required by the evolution of capitalism. It may be assumed that the more-than-one-party form constitutes a more accurate representation of the movement's reality than the one- party alternative. But the point made earlier about the potentially contradictory demands of representativity and effectiveness also applies here: a more-than-one-party situation in all likelihood produces greater representativity; but it may also make for less effectiveness. However, this is to say no more than that differences between people and groups make it more difficult than would otherwise be the case for them to take decisions. That is so. But there is a limit to what can be done about it, except by constraint and imposition.

Whether it makes for lesser political effectiveness or not—and the answer will vary according to different circumstances and factors—the more-than-one-party situation does not resolve the problem of 'substitutism', and it may not even greatly affect it one way or another. This is because two parties or more, purporting to represent the working class and the labour movement, may well take decisions that belong to the politics of leadership rather than the politics of the party as a whole. Indeed, it is likely that the more-than-one-party situation enhances the politics of leadership, in so far as it requires an often complex set of negotiations between allies or potential allies, and such negotiations emphatically form part of the politics of leadership.

So, from another angle, do most forms of revolutionary transformation, and notably revolutionary transformation by way of an insurrectionary seizure of power. A revolutionary movement launched on such a venture may well enjoy the support of a majority, even of an overwhelming majority, which may for instance be yearning to be freed from a hated regime. But this is by definition a matter of surmise; and revolutions, even in the best of circumstances, are not made by majorities, least of all revolutions which are made by way of insurrection. There may be no other way. But revolutions are made by minorities, and have usually been the work of relatively small minorities. At the very least, some element of 'substitutism' is here inevitable; and the refusal to engage in it may be fatal. Rosa Luxemburg is a good case in point. Her draft programme adopted by the founding congress of the German Communist Party in December 1918 included the following declaration: 'The Spartacus Union will never take over the power of government otherwise than by a clear manifestation of the unquestionable will of the great majority of the proletarian mass of Germany. It will only take over the power of government by the conscious approval of the mass of  the workers of the principles, aims, and tactics of the Spartacus Union.'\postnote{H. Gruber, International Communism in the Era of Lenin, p. 114 18 FI, p. 91.}

This indeed is a rejection of 'substitutism'. We do not know what would have happened in Germany if Luxemburg and her comrades, in a period of extreme political and social crisis, had viewed matters differently, and been less oppressed by the fear of usurping the popular will. But we do know that one month after she drafted the programme, she and Liebknecht were dead, and the old order was safe. Too great a propensity to 'substitutism' may turn into 'Blanquism' and lead to catastrophe. But so, in certain circumstances, may the rejection of 'substitutism' lead to defeat and catastrophe. In any case, Rosa Luxemburg's conditions, as set out in the above quotation, are most probably incapable of ever being met to the letter. In other words, revolutions are not only bound to include a certain element of 'substitutism' but actually to require it. This too weighs upon the Marxist project, and must at least to some extent affect the exercise of post-revolutionary power, and may affect it very greatly.

The discussion of class and party has so far proceeded as if the working class could only express itself politically by way of a party or parties. But this of course is not the case: although parties are its most important means of expression, there are other forms of working-class organization which have a direct bearing on political issues and struggles. One of these forms is trade unionism; another is 'conciliar' forms of organization— workers' councils, soviets, councils of action, and the like. These forms of organization also affect the whole question of the relation of class and party.

As far as trade unions are concerned, Marx and Engels saw them as performing a dual role, which they wanted to see linked. They believed that trade unions were the natural and indispensable product of the permanent struggle of labour against capital over all aspects of the work situation; and they also wanted them to be instruments in the political battle that had to be waged for the victory of labour over capital.

In his Instructions for Delegates to the Geneva Congress of the First International in September 1866, Marx expressed this duality by stating that 'if the trade unions are required for the guerilla fights between capital and labour, they are still more important as organized agencies for superseding the very system of wage labour and capital rule'.18 This formulation did no more than recognize a fact which was very familiar to Marx, given his knowledge of English trade unions, namely that unions might well seek to improve the conditions of labour within capitalism without being and without seeking to become 'organized agencies for superseding the very system of wage labour and capital rule'.

'Too exclusively bent upon the local and immediate struggles with capital', Marx also wrote in the same Instructions, 'the trade unions have not yet fully understood their power of acting against the system of wage slavery itself'; and he urged that apart from their original purposes, they must now learn to act deliberately as organizing centres of the working class in the broad interest of its complete emancipation. They must aid every social and political movement tending in that direction.'\postnote{Ibid., pp. 91-2.}

What is particularly notable about these and many other such formulations is not only that they show the crucial importance which Marx attached to the trade unions as political organizations; but also that he consistently failed to relegate the unions to the somewhat secondary and limited role, as compared to the party, which was assigned to them in later Marxist thinking. However this may be judged, it follows from the fact that neither Marx nor Engels had a particularly exalted view of 'the party' as the privileged expression of the political purposes and demands of the working class; and they did not therefore find it difficult to ascribe to trade unions a politically expressive role only a little less significant than that of 'the party'. This is not to suggest any kind of syndicalist streak in Marx, but only to note yet again that his emphasis on the working class itself and on the role which it must play in its own emancipation led him—whether rightly or wrongly—to be much less concerned than were later Marxists with the assignment of a towering role to the party as compared with that of the trade unions.\footnote{In a speech to a delegation of German trade unionists in 1869, Marx said that 'if they wish to accomplish their task, trade unions ought never to be attached to a political association or place themselves under its tutelage; to do so would e to deal them a mortal blow. Trade unions are the schools of socialism. It is in trade unions that workers educate themselves and become socialists because under their very eyes and every day the struggle with capital is taking place. Any political party, whatever its nature and without exception, can only hold the enthusiasm of the masses for a short time, momentarily; unions, on the other hand, lay hold on the masses in a more enduring way; they alone are capable of representing a true working-class party and opposing a bulwark to the power of capital.' (D. McLellan, The Thought of Karl Marx (London, 1971) p. 175).}

The relative devaluation of trade unions as political agencies—and it is no more than a relative devaluation—is quite marked in Lenin. I have already noted in a previous chapter Lenin's famous remark in What is to be Done? that 'the working class, exclusively by its own effort, is able to develop only trade-union consciousness'\postnote{V. I. Lenin, What is to be Done?, p. 375.} The designation is intended to suggest, as Lenin makes perfectly clear, a predominant concern of trade unions with immediate and limited economic demands associated with 'wages, hours and conditions'. This for Lenin did not mean that such concerns and demands would not have a political charge. On the contrary, he believed, as did Marx and Engels,\footnote{'Every class struggle is a political struggle' (Manifesto of the Communist Party, in Revs., p. 76).} that all struggles for economic improvements and reforms had and could not but have a certain political dimension; but also that the political aspects of such struggles would remain confined to its specific economic purposes. 'Trade- union consciousness' does not mean the absence of politics: it means the absence of revolutionary politics. For the latter to be present and for the transcendence of 'trade-union consciousness', what was required was an agency other than trade unions, namely the revolutionary party. Trade unions in the Leninist perspective are of very great importance; but nothing like as important, in terms of revolutionary politics, as the party.

One of the party's tasks is, in fact, to combat 'trade-union conciousness' in the trade unions themselves, and to treat them as arenas in which revolutionaries must struggle against trade- union leaders and a 'labour aristocracy' determined to prevent the transcendence of 'trade-union consciousness'. One of the sections of Lenin's 'Left-Wing' Communism - An Infantile Disorder was entitled 'Should Revolutionaries Work in Reactionary Trade Unions?', and his answer was an unqualified yes. 'We are waging a struggle against the “labour aristocracy'”, he wrote, 'in the name of the masses of the workers and in order to win them over to our side ... To refuse to work in the reactionary trade unions means leaving the insufficiently developed or backward masses of workers under the influence of the reactionary leaders, the agents of the bourgeoisie, the labour aristocrats, or “workers who have become completely bourgeois".'\footnote{SWL, p. 541. Lenin s quotation, he noted, was taken from a letter of Engels to Marx about British workers in 1858.}

In the same text, written, it will be recalled, in 1920, Lenin also noted that 'the trade unions were a tremendous step forward for the working class in the early days of capitalist development'; but that 'when the revolutionary party of the proletariat, the highest form of proletarian class organisation, began to take shape\dots the trade unions inevitably began to reveal certain reactionary features, a certain craft narrow-mindedness, a certain tendency to be non-political, a certain inertness, etc.' But he also went on to say that however, the development of the proletariat did not, and could not, proceed anywhere in the world otherwise than through the trade unions, through reciprocal action between them and the party of the working class'.\postnote{SWL, p. 539.} These formulations express clearly enough the basic Leninist view of trade unions as an indispensable element of class struggle, but also as an inevitably inadequate one, whose inadequacy must be remedied by the party. What this meant for the exercise of Bolshevik power is also spelt out very clearly by Lenin in Left-Wing Communism (and elsewhere), and will be considered presently. But it may be noted here that, while Lenin speaks of reciprocal action between trade unions and the party, he also very firmly asserts the primacy of the party, 'the highest form of proletarian class organisation'. In so far as the party is involved in some degree of 'substitutism', the trade unions can only help, at best, to reduce its extent, provided 'reciprocal action' can be given effective meaning.

The question of soviets, workers' councils, and 'conciliar' forms of organization in general, poses an even more direct challenge to Marxist political theory and practice in relation to class and party than does the question of trade unions. The challenge lies in the fact that councils are in one form or another a recurring and spontaneous manifestation of popular power in history; and that such forms of power come up, at least in twentieth-century revolutionary history, against the form of power represented by the workers' party or parties.

There is no guidance at all to be had here from Marx, since the key text from him on the issue of popular power, namely The Civil War in France, is entirely free of any reference to political parties of the working class. This text, written in the last weeks of the Paris Commune, is a glowing celebration of the form of popular power which he presented the Commune as having inaugurated or at least foreshadowed—'the political form at last discovered under which to work out the economical emancipation of labour'.\postnote{FI, p. 212.}
\begin{displayquote}
    Instead of deciding once in three or six years which member of the ruling class was to misrepresent the people in parliament, universal suffrage was to serve the people, constituted in communes, as individual suffrage serves every other employer in the search for the workmen and managers in his business. And it is well known that companies, like individuals, in matters of real business know how to put the right man in the right place, and, if they for once make a mistake, to redress it promptly.\postnote{Ibid., p. 210.}
\end{displayquote}

There is much else about this 'political form' which will need to be considered later, but the important point in the present context is that 'the people' is not, to all appearances, organized by anybody, nor is its relation with its representatives mediated, directed, or guided by a political party.

Nor, incidentally, did the passage of time cause Marx to place greater emphasis than he had done in The Civil War in France on the question of the lack of adequate organization and leadership from which the Commune obviously suffered. In a letter of 1881 he noted that the Commune 'was merely the rising of a city under exceptional conditions', and that its majority 'was in no wise socialist, nor could it be'.\postnote{SC, p. 410.} Obviously (and rightly), he held that conditions were not right for the Commune's success: 'With a modicum of common sense', he wrote in the same letter, 'it could have reached a compromise with Versailles useful to the whole mass of the people—the only thing that could be reached at the time.'\postnote{Ibid., p. 410.} But nowhere does he suggest that a different and more favourable outcome could have been achieved had the Communards been better organized, or even that any such venture imperatively requires organization, parties, leadership, etc. He did say in a letter written in April 1871, at the time when The Civil War in France was being composed, that the Central Committee of the National Cuard 'surrendered its power too soon, to make way for the Commune';\postnote{K. Marx to L. Kugelmann, 12 April 1971, in SC, p. 319.} and he would undoubtedly have welcomed a better organization of the Commune's endeavours. But it also appears that the absence of organization is not what mainly preoccupied him, or even that it greatly concerned him at all. This shows the dramatic shift of emphasis which occurred on this question in subsequent years.\footnote{It is instructive, in this context, to compare Marx with Trotsky, who expressed here a commonly accepted view: The Commune', he wrote, 'shows us the heroism of the working masses, their capacity to unite in a single bloc, their willingness to sacrifice themselves in the name of the future, but it also shows us at the same time the incapacity of the masses to choose their way, their indecision in the direction of the movement, their fatal inclination to stop after the first success, thus enabling the enemy to get hold of himself, to restore his position'. The Parisian proletariat', he went on, 'had no party, no leaders, to whom it would have been closely linked by previous struggles'; and again, 'it is only with the help of the party, which anticipates theoretically the paths of development and all its stages and extracts from this the formula of necessary action, that the proletariat frees itself from the necessity of ever beginning its history anew: its hesitations, its lack of decision, its errors' (Preface to C. Tales, La Commune de 1871 (Paris, 1921), pp. viii-ix).}

The idea of workers' concils and of 'communal' or 'conciliar' power formed no significant part of Marxist thought after the Paris Commune,\footnote{This is not to forget Engels's Introduction of 1891 to The Civil War in France but neither this nor the pamphlet itself affected the movement away from notions of 'conciliar' power.} until it suddenly forced itself upon the attention of Marxists in the form of the soviets in the Russian Revolution of 1905. In the light of later Bolshevik proclamations of their crucial importance, it may seem surprising that their emergence should have owed nothing to Marxist initiative, but it is nevertheless so.

In fact, there was much Bolshevik suspicion and even hostility to the Petersburg Soviet when it was set up in October 1905.\footnote{See e.g. M. Liebman, Leninism under Lenin, pp. 86 ff. The Petersburg Soviet represented 250,000 workers.} In essence, this attitude was a reflection of a dilemma which was at the heart of Bolshevik perspectives—on the one hand, the stress on the absolutely central role of the party in the revolutionary movement and of its guidance of the movement; on the other, the emergence of popular movements and institutions which owed nothing to the party and which could not be expected to be brought easily—or at all—under its guidance and control.

In 1905, Lenin responded differently from most Bolsheviks to the eruption of the Soviets on the political scene. In a letter to the Editor of the Party paper Novaya Zhizn (who did not publish it), Lenin, writing from Stockholm on his way back to Russia, denied that the question facing the Party was 'the Soviet of Workers' Deputies or the Party?'. On the contrary, he said, 'the decision must certainly be: both the Soviet of Workers' Deputies and the Party'; and he added that 'the only question—and a highly important one—is how to divide, and how to combine, the tasks of the Soviet and those of the Russian Social- Democratic Labour Party'.\postnote{V. I. Lenin, 'Our Tasks and the Soviet of Workers' Deputies', in CWL, vol. 10 (1962), p. 19.}

This was indeed the question. But the reflux of the revolutionary movement removed it from the Marxist agenda until 1917, when new Soviets again confronted the Bolsheviks with the same question, but this time with much greater insistence and urgency. It was then, in the months immediately preceding the October Revolution, that Lenin took up the question of the Soviets in his The State and Revolution. But it is precisely one of the most remarkable features of that extraordinary document that it not only fails to answer the question of the relation between the Soviets and the Party, but that it barely addresses itself to it.

The State and Revolution is an apotheosis of popular revolutionary power exercised through the Soviets of Workers' and Soldiers' Deputies. But these deputies would be subject to recall at any time, and they would operate within very strict limits, imposed upon them by the fact that they, in common with all other organs of power, would be subordinate to the proletariat. This and other aspects of The State and Revolution will be discussed later; but the crucial point here is that the party barely makes an appearance anywhere in the text. In fact, the only relevant reference to the party occurs in the following passage:
\begin{displayquote}
    By educating the workers' party, Marxism educates the vanguard of the proletariat, capable of assuming power and leading the whole people to socialism, of directing and organising the new system, of being the teacher, the guide, the leader of all the working and exploited people in organising their social life without the bourgeoisie and against the bourgeoisie.\postnote{The State and Revolution, in SWL, p. 281.}
\end{displayquote}

But even though the party only makes a fleeting appearance in The State and Revolution, there is no doubt that Lenin did intend 'the vanguard' to lead the people. In an article written even more closely to the time of the Bolshevik seizure of power, and in anticipation of it, he spoke glowingly of the Soviets as an entirely new form of state apparatus, immeasurably more democratic than any previous one, and he also noted that one of its advantages was that
\begin{displayquote}
    it provides an organisational form for the vanguard, i.e. for the most class-conscious, most energetic and most progressive section of the oppressed classes, the workers and peasants, and so constitutes an apparatus by means of which the vanguard of the oppressed classes can elevate, train, educate, and lead the entire vast mass of these classes, which has up to now stood completely outside of political life and history.\postnote{V. I. Lenin, Can the Rolsheviks Retain State Power?, ibid., p. 373. Notwithstanding its title, this was written on the eve of October.}
\end{displayquote}

What is missing from these and other such texts is any indication of how 'the vanguard', which must either be taken to be the party, or of which the party must be taken to be the core, would relate to those whom it wished to lead; and there is no indication either of what mechanisms would resolve the possible—indeed the inevitable—divergences between them. In other words, The State and Revolution did not answer what Lenin had in 1905 called 'the highly important' question of 'how to divide, and how to combine, the tasks of the Soviet and those of the Russian Social-Democratic Party'. The reason for that failure was that Lenin in 1917 was carried ahead by an extreme revolutionary optimism, fostered by the advance of the popular movement, and which led him to ignore the question he had posed in 1905. His optimism was of course unwarranted; and the Bolsheviks were soon compelled by the circumstances which they confronted to answer the question in very different ways from those which Lenin had envisaged when he hailed the 'new state apparatus' which he believed to be foreshadowed by the Soviets.

\section*{3}
The State and Revolution is the most authoritative text in Marxist political writing on the dictatorship of the proletariat a concept to which Marx attached supreme importance but which he never defined in any detail.\footnote{In his Introduction of 1891 to The Civil War in France, Engels, addressing the Social-Democratic philistine' who was 'filled with wholesome terror at the words: Dictatorship of the Proletariat', had said: 'Do you want to know what this dictatorship looks like? Look at the Paris Commune. That was the Dictatorship of the Proletariat' (SW 1968, p. 262). Marx himself never said this, though it is reasonable to assume that his description of the Commune's organization was a close approximation to the main features of the dictatorship of the proletariat, as he conceived it.} In the Critique of the Gotha Programme (1875), Marx had written that 'between capitalist and communist society lies a period of revolutionary transformation from one to the other. There is a corresponding period of transition in the political sphere and in this period the state can only take the form of a revolutionary dictatorship of the proletariat.'\postnote{FI, p. 355.} Lenin quite rightly emphasized again and again that its basic premise was, as Marx had said in The Civil War in France, that 'the working class cannot simply lay hold of the ready-made state machinery, and wield it for its own purposes'.\footnote{FI, p. 206. He and Engels thought the point sufficiently important to single out the quotation and reproduce it in their 1872 Preface to the German Edition of the Communist Manifesto (SW 1968, p. 32); and Engels reproduced it once again in his Preface to the English Edition of 1888 (SW 1950, I, p. 29).} In the letter to Kugelmann to which I have already referred, Marx also reminded his correspondent that in the Eighteenth Brumaire, T declare that the next attempt of the French Revolution will be no longer, as before, to transfer the bureaucratic\hyp{}military machine from one hand to another, but to smash it, and this is the preliminary condition for every real people's revolution on the Continent. And this is what our heroic Party comrades in Paris are attempting'.\footnote{SC, p. 318. In the Eighteenth Brumaire, Marx had said that 'all political upheavals perfected this machine instead of smashing it' (K. Marx.SE, p. 238).} It was this 'smashing' of the existing state which Lenin, following Marx, took as the first and absolutely essential task of a genuinely revolutionary movement and party.

What would follow, he explained, was 'a gigantic replacement of certain institutions by other institutions of a fundamentally different type'.\postnote{V. I. Lenin, The State and Revolution, SWL, p. 293.} Here too, Lenin was faithfully following Marx's pronouncements on the Commune in The Civil War in France. 'The first decree of the Commune', Marx had said, 'was the suppression of the standing army, and the substitution for it of the armed people';\postnote{FI, p. 209.} and he then went on to list other major features of the Commune's organization:
\begin{displayquote}
    The Commune was formed of the municipal councillors, chosen by universal suffrage in the various wards of the town, responsible and revocable at short terms. The majority of its members were naturally working men, or acknowledged representatives of the working class. The Commune was to be a working, not a parliamentary body, executive and legislative at the same time. Instead of continuing to be the agent of the central government, the police was at once stripped of its political attributes, and turned into the responsible and at all times revocable agent of the Commune. So were the officials of all other branches of the administration. From the members of the Commune downwards, the public service had to be done at workmen's wages. The vested interests and the representation allowances of the high dignitaries of state disappeared along with the high dignitaries themselves.\postnote{Ibid., p. 209.}
\end{displayquote}

It was this which for Lenin constituted a whole programme, precisely the 'gigantic replacement of certain institutions by other institutions of a fundamentally different type', these other institutions being the Soviets, which provided 'a bond with the people, with the majority of the people, so intimate, so indissoluble, so easily verifiable and renewable, that nothing remotely like it existed in the previous state apparatus'.\postnote{Can the Bolsheviks Retain State Power?, SWL, p. 373.} Moreover, the Soviets made it possible 'to combine the advantages of the parliamentary system with those of immediate and direct democ- racy, i.e. to vest in the people's elected representatives both legislative and executive functions'.\postnote{Ibid., p. 374.}

There can be no doubt that Lenin believed that the Soviet system, once in operation, would represent the most democratic and popular regime that could be achieved before the advent of a fully socialist society. Compared with the bourgeois parliamentary system', he also wrote, 'this is an advance in democracy's development which is of world-wide, historic significance'.\postnote{Ibid., p. 374.} A few weeks after this was written, the Bolsheviks had taken power, and the Soviet system continued to spread all over Russia. At the Third All-Russia Congress of Soviets which assembled in January 1918, Lenin felt able to say to the delegates: 'Look wherever there are working people, look among the masses, and you will see organizational, creative work in full swing, you will see the stir of a life that is being renewed and hallowed by the revolution.'\postnote{Liebman, op. cit., p. 219.}

This conviction that a new political order of an unprecedented popular and democratic kind was actually being born was obviously not the only or even the main reason for the momentous decision which Lenin and the other Bolshevik leaders took to dissolve the Constituent Assembly as soon as it had met at the beginning of January 1918. But the belief that there was a genuinely Soviet-type alternative to the Constituent Assembly undoubtedly helped to reinforce the political and ideological considerations which led to that decision.\footnote{The most important of these considerations was obviously the fact that the the Constituent Assembly, which had been held in November 1917, had returned a majority of opponents of the Bolsheviks. Whatever view is taken of the Assembly's dissolution, the Bolsheviks were obviously right to think that their rule was incompatible with the Assembly's continued existence.} Lenin did speak in the first months after the October Revolution as if he believed that it had inaugurated a regime of popular and democratic power such as he had outlined in The State and Revolution, in which the problem of class and party would be resolved in a sort of symbiotic relation between them. On this basis, it was obviously possible to think not only that the Constituent Assembly would place a dangerous brake—at the least—upon a still developing revolutionary movement but that, from the point of view of democracy itself, it was also an unnecessary encumbrance and a regression to parliamentarism at a time when new and immensely more democratic forms of political power had come into being and were spreading.

In fact, the faith which Lenin had placed in the Soviets was rendered altogether illusory by the circumstances of revolution and civil war. Whether they could have fulfilled even some of his expectations had circumstances been more favourable is an open question. But even if they had, the relation between the Party and the Soviets would still have posed the question which Lenin had asked in 1905. As it was, the disintegration of the Soviets and the terrible dangers to which the revolution was exposed from within and from without vastly favoured the 'substitutism' which some revolutionaries had previously feared and which some in the Bolshevik ranks virulently attacked after 1917.

Lenin responded to such attacks by denying that the party was usurping the place of the proletariat and by pleading exceptional conditions. His change of perspective from the period preceding and following the Bolshevik seizure of power is well indicated by his claim in 1919 that 'the dictatorship of the working class is carried into effect by the party of the Bolsheviks which since 1905 has been united with the whole revolutionary proletariat'.\postnote{E. H. Carr, The Bolshevik Revolution, 1917-1923 (London, 1950) p. 230.} It was also the same view of the special character of the Bolshevik Party which led him defiantly to accept the charge of 'dictatorship of one party': 'Yes, the dictatorship of one party! We stand upon it and cannot depart from this ground, since this is the party which in the course of decades has won for itself the position of vanguard of the whole factory and industrial proletariat.'\postnote{Ibid., p. 230.} 

It is true that Lenin and the Bolsheviks were alone capable of saving what was after ail their revolution; and there was no real possibility of alliance with any other group or party.\postnote{For Lenin's attitude to other parties and groups, and vice versa, see Liebman, op. cit., pp. 238 ff.} But what this situation entailed was the dictatorship of the Bolshevik Party over the proletariat as well as over all other classes, by means of the repressive power of the state which the party controlled. Lenin might well deny that it was appropriate to speak of a stark dichotomy between "dictatorship of the party or dictatorship of the class',\postnote{See the reference to 'Left-Wing' Communism—An Infantile Disorder, p. 123 above.} but the fact remained that the party had 'substituted' itself for a ravaged and exhausted working class, in a country gripped by civil war, foreign intervention and economic collapse.

But the growing concentration of power which these circumstances encouraged very quickly came to affect the party itself. E. H. Carr notes that 'the seventh party congress of March 1918 which voted for the ratification of Brest-Litovsk was the last to decide a vital issue of policy by a majority vote\dots Already in October 1917 it was the central committee which had taken the vital decision to seize power; and it was the central committee which succeeded to the authority of the congress.'\postnote{Carr, op. cit., p. 193.} But this too, as he also notes, was soon followed by the exercise of power by more restricted and secretive organs.

The changes that were occurring in the nature and spirit of the party found formal expression in the decision taken at the Tenth Party Congress in March 1921 at the instigation of Lenin to forbid all groups and factions within the party and to give to the Central Committee full powers to punish, if need be by expulsion, 'any breach of discipline or revival or toleration of fractionalisin'.\postnote{Ibid., p. 201.} The previous three years, although full of crises and dangers, had been marked by considerable controversy in the ranks of Bolshevism, or at least among its leading personnel.

The polemics had ranged over many issues, notably over the proper application of 'democratic centralism' and the role of the trade unions in the new Soviet state.\footnote{The first major issue had been whether to Litovsk. accept the peace terms of Brest-Litovsk.} In the debate on the trade unions one faction, led by Trotsky, wanted the thorough subordination of the unions to the state; another, the 'workers' opposition', wanted much greater independence for them and their assumption of far larger powers in the running of economic and  industrial life.\footnote{For the debate on the trade unions, see e.g. I. Deutscher. Soviet Trade Unions (Their Place in Labour Soviet Policy) (London, 1952).} Lenin occupied a middle position, and rejected both the 'militarization' of the unions which Trotsky advocated, and also an independence for them that would have threatened the new regime with a 'syndicalist' deformation. But by the beginning of 1921, the 'Workers' Opposition' had emerged as a distinctive group, with a substantial measure of support, and with a programme sufficiently loose to attract a great deal more support at a time of acute economic stress and political crisis (the Tenth Party Congress met as the Kronstadt rising was being crushed). It was in these circumstances that the decision was ratified by that Congress to enforce 'unity' upon the party.

There is no question that it was Lenin who bore the prime responsibility for this major step forward in the direction of what 1 have referred to earlier as 'dictatorial centralism'. But it is of interest that he was by no means the most extreme advocate of that decision and that he did not try to make a virtue of what he believed to be required by dire circumstances ('We need no opposition, comrades, now is not the time').\postnote{Ibid., p. 199.} Ironically, Trotsky who had so virulently denounced Lenin's 'substitutism' in 1903 unequivocally claimed in 1921 the party's right 'to assert its dictatorship even if that dictatorship temporarily clashed with the passing moods of the workers' democracy\dots The party is obliged to maintain its dictatorship, regardless of temporary wavering in the spontaneous moods of the masses, regardless of the temporary vacillations even in the working class.'\footnote{Deutscher, The Prophet Armed, op. cit., p. 509. Note, however, Trotsky's further remark that the dictatorship does not base itself at every given moment on the formal principle of a workers' democracy, although the workers' democracy is, of course, the only method by which the masses can be drawn more and more into the political life' (ibid.). The strain in thought is here evident.}

However, it was also Trotsky who, in a new Preface to his book 1905, dated 12 January 1922, came closest to the heart of the problem which had confronted the Bolsheviks. Referring to the period from March to October 1917, he wrote that 'although having inscribed on our banner: “All power to the Soviets”, we were still formally supporting the slogans of democracy, unable as yet to give the masses (or even ourselves) a definite answer as to what would happen if the cogs of the wheels of formal democracy failed to mesh with the cogs of the Soviet system'; and he then went on:
\begin{displayquote}
    The dispersal of the Constituent Assembly was a crudely revolutionary fulfilment of an aim which might also have been reached by means of a postponement or by the preparation of elections. But it was precisely this peremptory attitude towards the legalistic aspects of the means of struggle that made the problem of revolutionary power inescapably acute, and in its turn, the dispersal of the Constituent Assembly by the armed forces of the proletariat necessitated a complete reconsideration of the interrelationship between democracy and dictatorship.\footnote{L. Trotsky 1905 (London, 1971), pp. 21-2. Trotsky added that 'In the final analysis, this represented both a theoretical and a practical gain for the Workers' International', which is obviously absurd.}
\end{displayquote}

It was precisely this complete reconsideration' which the Russian Communists were unable to undertake in Lenin's lifetime, and which they were forbidden to undertake after his death. By the Twelfth Party Congress in 1923, which Lenin's illness prevented him from attending, Zinoviev was already proclaiming that we need a single strong, powerful central committee which is leader of everything',\postnote{Ibid., p. 231.} and it was, of all people, Stalin who demurred, since he did not want power concentrated in the Central Committee at the expense of the Secretariat and other working organs of the Party.\postnote{Ibid., p. 231.}

This evolution from October 1917 until 1923 - the years of Leninism in power has inevitably been considered, particularly since 1956, in the light of the grim experience of Stalinism; and it has become a familiar argument that Leninism 'paved the way for Stalinism, and that the latter was only the 'natural' successor of the former.

The argument turns on the meaning which is given to Stalinism. If it is merely taken to mean a further centralization of power than had already occurred and a greater use of the repressive and arbitrary power than was already the case in the Leninist years, then the continuity is indeed established.

But Stalinism had characteristics and dimensions which were lacking under Leninism, and which turned it into a regime much more accurately seen as marking a dramatic break with anything that had gone before. The most important of these must be briefly noted here.

To begin with, it was a regime in which one man did hold absolute power of a kind which Lenin never remotely had (or ever gave the slightest sign of wanting to achieve); and Stalin used that power to the full, not least for the herding into camps of millions upon millions of people and the 'liquidation' of countless others, including vast numbers of people who were part of the upper and upmost layers of Soviet society—this devastation of all ranks of the regime's officialdom, which was a recurring phenomenon, is unique as an historical event. The sheer scale of the repression is a second feature of Stalinism which distinguishes it most sharply from Leninism. The Soviet Marxist (and 'dissident') historian Roy Medvedev put the point in a striking formulation by saying that 'the NKVD arrested and killed, within two years, more Communists than had been lost in all the years of the underground struggle, the three revolutions and the Civil War';\postnote{R. Medvedev, Let History? fudge: The Origins and Consequences of Stalinism (London, 1972), p. 234.} and he also quotes figures which suggest that for three years of the Civil War, 1918—20, fewer than 13,000 people were shot by the Cheka.\postnote{Ibid., p. 390.} Even if the figures are contested as being far too low, the difference in scale between the repression that occurred in the years of revolution, civil war, and foreign intervention on the one hand, and the holocaust of Stalinism on the other, is so vast as to discredit the notion of continuity. Moreover, repression on such a scale required a gigantic police apparatus that reached out to every corner of Soviet life, not only to arrest, deport, and execute, but to maintain detailed surveillance over Soviet citizens who were not imprisoned and to guard those who were. This was an exceedingly elaborate system which was quite deliberately woven into the tissue of the larger social system and that decisively affected the latter's total pattern.

This points to a third characteristic of Stalinism, namely that it required from the people a positive and even enthusiastic acceptance of whatever Tine', policy, position and attitude was dictated from on high, however much it might contradict the immediately preceding one. This applied not only to major aspects of policy, or even to minor ones, but to every conceivable aspect of life and thought. Conformity was demanded over literature and music as well as over foreign policy and Five-Year plans, not to speak of the interpretation of Marxism-Leninism and of history, particularly the history of the Russian Revolution; and every vestige of opposition was ruthlessly stamped out. Indeed, Stalinism was a regime which stamped out opposition in anticipation, and constantly struck at people who were perfectly willing to conform, on suspicion that they might eventually cease to be willing. No such conformity was demanded in the first years of the revolution; and it formed no part of Leninism.

This Stalinist requirement of total conformity to all Soviet policies and actions was also extended to the world Communist movement. Whenever the regime could, it physically stamped out opposition or suspected opposition among foreign Communists; and it demanded the same rigid endorsement of every single aspect of Russian internal and external policy that it was able to exact inside the Soviet Union. (That it recieved this ensorsement is of course one of the most remarkable and still very poorly charted phenomena of the twentieth century.) What happened was the total Stalinization of every single Communist Party throughout the world, in the name of the sacred duty imposed upon every Communist to defend the U.S.S.R.; and defending the U.S.S.R. rapidly came to be interpreted as including the defence of every twist and turn in Soviet internal as well as external policy; the endorsement of whatever Stalinist 'line' might be current on literature, biology, linguistics, or whatever; the denunciation as traitors, renegades, and foreign agents of anyone who was so labelled by the Soviet leadership, including all but a few of those who had been the leading figures of Bolshevism; and the complete and vehement rejection as bourgeois (or Fascist) lies of any charge levelled at the rulers of the socialist sixth of the world'. This kind of automatic submission to Russian requirements formed no part of Leninism either.

In the perspective of the present chapter, another characteristic of Stalinism must be stressed: this is the apotheosis of 'the Party'. In so far as 'the Party' meant in fact its leadership, and above all its supreme leader, this could be taken as part of the cult of personality . But the distinction between the latter and what might be called the cult of the party nevertheless needs to be made. For it became habitual and indeed necessary after Lenin s death to speak of 'the Party ' in quasi-religious terms and to invoke its unity as a quality that must on no account be infringed, and whose infringement warranted the direst penalties. This was not only of the utmost value to those who controlled the Party: it was also a paralysing constraint on those who opposed them. As early as the Ihirteenth Party Congress in May 1924, the first after Lenin s death, Trotsky, who had belatedly identified himself with the cause of inner-party democracy, nevertheless still spoke of the Party in terms which were of  greater comfort to his enemies in power than to his allies in opposition: 'Comrades', he said, 'none of us wishes to be or can be right against the party. In the last instance the party is always right, because it is the only historic instrument which the working class possesses for the solution of its fundamental tasks'.\postnote{I. Deutscher, The Prophet Unarmed. Trotsky: 1921-1929 (London, 1959), p. 139.}

He refused to recant, as he had been required to do by Zinoviev, who was then allied to Stalin; and he said that he did not believe that all the criticisms and warnings of those who opposed the leadership were wrong 'from beginning to end'. But he also said 'I know that one ought not to be right against the party. One can be right only with the party and through the party because history has created no other way for the realization of one's rightness.'\postnote{Ibid., p. 139.}

Trotksy soon changed his mind. But in one form or another, such debilitating notions continued to help disarm generation after generation of revolutionaries in all Communist parties, and must be taken as part of the explanation for the extreme weakness of the Marxist opposition to Stalinism in subsequent years. A conviction was deeply implanted—and Stalinist propaganda did all it could to strengthen it—that outside the party, there could be no effective action; and inside the party, there was total control from the top. This was the case for all Communist parties. Opponents in these parties were either subdued or expelled, and in a number of cases physically destroyed.

Lenin had never made a cult of the party. He had in the stress of crisis identified the dictatorship of the proletariat with the dictatorship of the party; and his name is indissolubly linked with the idea that no revolutionary movement is conceivable without a revolutionary party to lead it. But he never invested 'the Party' with the kind of quasi-mystical attributes that became habitual after his death, and which had hitherto been more properly associated with religious than with secular institutions—but then a religious frame of mind did come to dominate the world Communist movement in the years of Stalinism, with 'Marxism-Leninism' as its catechism, 'dialectical materialism' as its mystery, 'the Party' as its Church, and Stalin as its prophet.

This frame of mind also made it easy to suppress the question of 'substitutism'; and it was indeed one of the hallmarks of Stalinism to turn this into a non-question. It was taken for granted that 'the Party' and the working class formed a perfect unity, and that the former by definition 'represented' the latter; and since there was no problem of 'substitutism', neither was there any need to look seriously for remedies or answers to it.

On the other hand, Lenin's anguished awareness in his last years that the problem was real and growing did not lead him to propose any adequate solution to it, either in the short term or the long.\postnote{See, e.g., M. Lewin, Lenin's Last Struggle (London, 1969).} In truth, he had no such solution to offer. Nor was there one to be had within the confines of a system that combined the monopoly of one party with dictatorial centralism.

The difficulties presented by the relation of party to class were enormously compounded by the specific and desperately unfavourable conditions of post-revolutionary Russia: but they did not stem from these conditions. The point is underlined by reference to the attempt which Gramsci made, quite independently of Russian experience or circumstances, to resolve the question theoretically, and by the fact that he was only able to restate it.

In an article, The Problem of Power', in L'Ordine Nuove, which appeared in November 1919, Gramsci wrote that this was
\begin{displayquote}
    the problem of the modes and forms which will make it possible to organise the whole mass of Italian workers in a hierarchy which organi- ca y culminates in the party. It is the problem of the construction of a state apparatus which, internally, will function democratically, that is, will guarantee freedom to all anti-capitalist tendencies, the possibility of becoming parties of proletarian government, but externally, will be like an implacable machine which crushes the organisations of the industrial and political power of capitalism.\postnote{L Ordine Nuove, 29 November 1919, in New Edinburgh Review. Gramsci, II, p. 73. The article is one of a number reproduced in the Review in a Section entitled 'Proletarian Forms' and edited by Gwyn Williams.}
\end{displayquote}

Commenting on this, Gwyn Williams suggests that 'the relationship between party and masses in this, as in other projects of a similar temper (Rosa Luxemburg's concept of leadership, for example) reads like an attempt to square the circle'.\postnote{Ibid., p. 73.} Whether so or not, it is clear that these formulations do not point to any solution of the question. The same has to be said of Gramsci's later Prison Notebooks. These include some illuminating reflections on the proletarian party's role as the 'modern Prince', the collective intellectual' of the working class and as its essential form of political organization. But they do not provide the theoretical material for the resolution of the tension between class and party. Nor, as I have suggested earlier, is it likely that such a resolution is wholly possible: what can be achieved is the attenuation of this tension; and the degree to which it is attenuated is a matter of crucial importance. But this requires  severe and effective limitations on the powers of leaders, which is no small undertaking and which is itself dependent on the existence or the achievement of a cluster of favourable circumstances.

The Chinese experience, which some people on the left have tended to regard as a demonstration of how the problem can be solved, is in fact an instructive example of its gravity. Not very paradoxically, this is so precisely because no Communist leader has been more explicitly and more consistently concerned than Chairman Mao with the relation of leaders to led and with the meaningful application of the principle of 'democratic centralism': yet the Chinese record, in this realm at least, is far from impressive. It may well be argued that it is better than the Russian one; but this is not saying much.

One of the most distinctive features of Maoism has been its constant stress on the danger that leaders and cadres would become isolated, rigid, and remote from the masses; and a related distinctive feature has been the conviction that the remedy lay not only in the acceptance but in the encouragement and fostering of criticism and challenge 'from below'. The following quotation from a talk by Chairman Mao in 1962 to a Conference of 7,000 cadres from various levels may serve as a typical instance of his approach to this issue:
\begin{displayquote}

    If there is no democracy, if ideas are not coming from the masses, it is impossible to establish a good line, good general and specific policies and methods. Our leading organs merely play the role of a processing plant in the establishment of a good line and good general and specific policies and methods. Everyone knows that if a factory has no raw material it cannot do any processing. If the raw material is not adequate in quantity and quality it cannot produce good finished products. Without democracy, you have no understanding of what is happening down below; the situation will be unclear; you will be unable to collect sufficient opinions from all sides; there can be no communication between top and bottom; top-level organs of leadership will depend on one-sided and incorrect material to decide issues, thus you will find it difficult to avoid being subjectivist; it will be impossible to achieve unity of understanding and unity of action, and impossible to achieve true centralism.\postnote{Mao Tse-tung Unrehearsed. Talks and Letters: 1956-71 (London, 1974), p. 164.}
\end{displayquote}

The first thing to note about this text is that the argument for 'democracy' is primarily a 'functional' one: cadres will be better informed and will reach better decisions if they listen and pay heed to what the masses are saying. What the masses are saying is the raw material which cadres have to process. The stress is always on 'letting the masses speak out' in order to improve the work of the Party. But there is no question that the Party must always lemain in command and be the final arbiter of what is to be done.

Chairman Mao was willing to go very far in trying to compel the Party to 'listen to the masses' and to 'apply the mass line — how far is best shown by the Cultural Revolution which was directed against 'persons in authority who had taken the capitalist road', and most of which were obviously Party cadres and Party members. One of the slogans put forward at the beginning of the Cultural Revolution, under the inspiration of the Chairman, namely 'Bombard Party Headquarters', is indicative of the spirit of challenge which was being generated at the time. (The argument, incidentally, that the whole operation was designed by the Chairman and his acolytes to get rid of his opponents or at least to subdue them, is of no great relevance here—the methods employed are what matters.)

On the other hand, it is equally significant that the Chairman never gave any indication that he had the slightest intention of surrendering or even reducing the 'leading role' of the Party; or o bringing about major institutional reforms in its mode of being. Whenever the movement that he had generated or encouraged gave any sign of really getting out of hand and of engulfing the Party, he and his partisans firmly reasserted the crucial importance of maintaining its directing character. Other organizations—the People's Liberation Army and revolutionary committees—could and did play a major role. But it was never suggested that they or any other organization could conceivably replace the Party or assume a co-equal role with it. Nor indeed was this conceivable in the circumstances of the Chinese Revolution.

It is rather more interesting that Chairman Mao never gave any sign that he was concerned to give an institutional and solidly constructed basis to the process of involvement of the masses which he consistently advocated; and that he never suggested at such involvement required the creation of institutions independent from Party control.

Popular involvement is not the same thing as democratic participation and control; and it is perfectly possible to have a very large measure of the one without having much (or even any) of the other. Millions upon millions of people have been involved in great and tumultuous movements in China; and their involvement has included the (officially encouraged) criticism of cadres and 'persons in authority'. Many of these have been swept from their offices and positions, either temporarily or for good. But however fruitful such movements of criticism and challenge may be judged to be, they are by no means the same as democratic participation and popular control and do not, in any serious meaning of the word, have much to do with either.

The kind of 'purges' which have been part of Maoism, notably in the Cultural Revolution, have generally speaking been a great improvement on the 'purges' of Stalinism: but they have not changed the system which made the 'purges' necessary in the first place. Chairman Mao spoke of the need to have a whole series of Cultural Revolutions stretching into the indefinite future. The trouble is that, in so far as such Revolutions do not greatly affect structures and institutions, they are not an adequate means of dealing with the degenerative processes with which the Chairman was rightly concerned.

Much may be claimed for the Chinese experience. But what cannot be claimed for it, on the evidence, is that it has really begun to create the institutional basis for the kind of socialist democracy that would effectively reduce the distance between those who determine policy and those on whose behalf it is determined. Nor for that matter has Maoism made any notable contribution to the theoretical attack of the question.

In terms of Marxist politics in general, the question remains open; and its discussion is best pursued by considering the Marxist view of the revolutionary process, with which the discussion of socialist democracy is intertwined.

\printpostnotes